\chapter{Random Matrices II: Unitary Matrices and the Zeros of Zeta}\chaptermark{Unitary Matrices and the Zeros of Zeta}

\begin{center}\itshape
Choose a unitary matrix at random and look at its eigenvalues on the circle.\\
Then compute the zeros of the Riemann zeta function, high up the critical line.\\
You will not be able to tell the two pictures apart, and nobody knows why.
\end{center}

In 1972 Hugh Montgomery, a young American number theorist visiting the Institute for Advanced Study in Princeton, had a formula in his pocket. Assuming the Riemann hypothesis, he had studied how the
zeros $\frac12+i\gamma$ of the zeta function are spaced, and he had found that the pairs of zeros seem to obey a law: the density of pairs at distance $u$ (in units of the mean spacing) should be $1-\bigl(\frac{\sin\pi u}{\pi u}\bigr)^2$.
The function vanishes at $u=0$: the zeros repel each other. At afternoon tea Sarvadaman Chowla introduced him to Freeman Dyson. Montgomery started to explain what he was doing, and Dyson, as Montgomery later told the
writer Karl Sabbagh, recognised the formula at once: it was the pair correlation of the eigenvalues of a random Hermitian matrix, which Dyson had computed ten years earlier while thinking about heavy nuclei.
A number theorist and a physicist had found the same function for reasons that had, apparently, nothing to do with each other.

Fifteen years later Andrew Odlyzko, at Bell Laboratories, computed $10^5$ consecutive zeros around the $10^{12}$-th zero, rescaled them, and drew the histogram of their spacings next to the random matrix prediction. The curves are
indistinguishable to the eye. In 2000 Jon Keating and Nina Snaith went one step further: they computed the moments of the characteristic polynomial of a random unitary matrix, and the result predicted the moments of the
zeta function on the critical line, including two constants, $42$ and $24024$, that number theorists were obtaining at the same time by much harder heuristics.

This chapter proves everything on the random matrix side. The eigenvalues of a random unitary matrix, seen at the scale of their spacing, form the \emph{sine process}; they are so rigid that the number of them in an arc
fluctuates only by $\sqrt{\log}$ of its length; and the characteristic polynomial factors, in law, into independent pieces, which gives its moments in closed form and a central limit theorem for its logarithm. The zeta side is
stated, and checked on ten thousand zeros computed numerically.

\begin{reminder}[from Volume 8, Volume~2, Volume 3 and Volume 10]
$U$ is Haar distributed on $U(n)$, with eigenvalues $e^{i\theta_1},\dots,e^{i\theta_n}$, $\theta_j\in[0,2\pi)$. Haar measure is left and right invariant and unique (Volume~8, Theorem~8.3). \emph{Heine's identity}: for $f$ bounded and
Borel on the circle, $\mathbb E\prod_jf(e^{i\theta_j})=\det\bigl(\hat f(j-k)\bigr)_{0\le j,k\le n-1}$ (Volume~8, Theorem~8.13). With $K_n(\theta,\varphi)=\sum_{a=0}^{n-1}e^{ia(\theta-\varphi)}$, the \emph{correlation functions} are determinants: for bounded Borel
$g$ on $[0,2\pi)^k$,
\[
\mathbb E\sum_{j_1,\dots,j_k\text{ distinct}}g(\theta_{j_1},\dots,\theta_{j_k})=\int_{[0,2\pi)^k}g(\varphi)\det\bigl(K_n(\varphi_a,\varphi_b)\bigr)_{a,b\le k}\frac{\dd\varphi_1\cdots\dd\varphi_k}{(2\pi)^k}
\]
(Volume~8, Theorem~8.15), and $|K_n(\theta,\varphi)|=\bigl|\frac{\sin(n(\theta-\varphi)/2)}{\sin((\theta-\varphi)/2)}\bigr|\le n$. Dyson's pair correlation $1-(\frac{\sin\pi u}{\pi u})^2$ is Volume~8, Corollary~8.16. \emph{Hadamard's inequality}: a positive semidefinite
matrix has determinant at most the product of its diagonal entries (Volume~2, Theorem~K.5). \emph{Lyapunov's condition}: if the rows $X_{n,1},\dots,X_{n,k_n}$ are independent and centred with $\sum_k\mathbb EX_{n,k}^2=1$ and
$\sum_k\mathbb E|X_{n,k}|^{2+\delta}\to0$, then $\sum_kX_{n,k}\Rightarrow N(0,1)$ (Volume~10, Corollary~4.2). Slutsky's lemma: Volume~3, Exercise~6.1.
\end{reminder}

\section{Counting eigenvalues}

For a Borel set $D\subset[0,2\pi)$ let $N(D)=\#\{j:\theta_j\in D\}$ and $|D|$ its length. The first surprise is that $N(D)$, a count of strongly dependent points, is a sum of \emph{independent} coin flips.

\begin{proposition}[a sum of independent Bernoulli variables]\label{prop:bernoulli}
Let $M_D$ be the $n\times n$ matrix $\bigl(\int_De^{i(a-b)\theta}\frac{\dd\theta}{2\pi}\bigr)_{0\le a,b<n}$. It is Hermitian with eigenvalues $p_1,\dots,p_n\in[0,1]$, and $N(D)$ has the law of $B_1+\dots+B_n$ with independent $B_j\sim\mathrm{Bernoulli}(p_j)$. In particular
\[
\mathbb EN(D)=\frac{n|D|}{2\pi},\qquad\operatorname{Var}N(D)=\operatorname{tr}M_D-\operatorname{tr}M_D^2.
\]
\end{proposition}

\begin{proof}
For $v\in\mathbb C^n$,
\[
v^*M_Dv=\sum_{a,b}\bar v_av_b\int_De^{i(a-b)\theta}\frac{\dd\theta}{2\pi}=\int_D\Bigl|\sum_a\bar v_ae^{ia\theta}\Bigr|^2\frac{\dd\theta}{2\pi}.
\]
This is at least $0$, and at most the same integral over the whole circle, which is $\|v\|^2$ by orthogonality of
the $e^{ia\theta}$. So $M_D$ is Hermitian with $0\le M_D\le I$, and its eigenvalues lie in $[0,1]$.

Fix $z\in\mathbb C$ and let $f=1+(z-1)\mathbf 1_D$, so that $\prod_jf(e^{i\theta_j})=z^{N(D)}$. Its Fourier coefficients are $\hat f(m)=\delta_{m0}+(z-1)\int_De^{-im\theta}\frac{\dd\theta}{2\pi}$, so the Toeplitz matrix of Heine's identity is
$I+(z-1)M_D^{\mathsf T}$, and
\[
\mathbb Ez^{N(D)}=\det\bigl(I+(z-1)M_D^{\mathsf T}\bigr)=\prod_{j=1}^n\bigl(1+(z-1)p_j\bigr)=\prod_{j=1}^n(1-p_j+p_jz).
\]
The right side is the generating function of $B_1+\dots+B_n$. Two random variables with values in $\{0,\dots,n\}$ whose generating polynomials agree have the same law. The mean and variance of the sum are
$\sum p_j=\operatorname{tr}M_D=\frac{n|D|}{2\pi}$ and $\sum p_j(1-p_j)=\operatorname{tr}M_D-\operatorname{tr}M_D^2$.
\end{proof}

The representation is an instance of a theorem of Ben Hough, Manjunath Krishnapur, Yuval Peres and B\'alint Vir\'ag (2006) on all determinantal point processes. It reduces every question about $N(D)$ to the eigenvalues of
$M_D$, and the variance needs only their sum and the sum of their squares.

\begin{theorem}[the number variance]\label{thm:numbervariance}
Let $D$ be an arc of length $2\pi\alpha$, and $L=n\alpha$ the expected number of eigenvalues in it. Then
\begin{equation}\label{eq:numbervariance}
\begin{split}
\operatorname{Var}N(D)&=\frac2{\pi^2}\sum_{m=1}^{n-1}\frac{\sin^2\pi m\alpha}m\\
&\quad+\frac{2n}{\pi^2}\sum_{m\ge n}\frac{\sin^2\pi m\alpha}{m^2},
\end{split}
\end{equation}
and there is an absolute constant $C$ such that $\bigl|\operatorname{Var}N(D)-\frac1{\pi^2}\log L\bigr|\le C$ whenever $1\le L\le\frac n2$.
\end{theorem}

\begin{proof}
By rotation invariance ($e^{i\varphi}U$ is again Haar distributed) we may take $D=[0,2\pi\alpha)$. Then $(M_D)_{ab}=c_{a-b}$ with $c_m=\int_0^{2\pi\alpha}e^{im\theta}\frac{\dd\theta}{2\pi}$, so $c_0=\alpha$ and $|c_m|=\frac{|\sin\pi m\alpha|}{\pi|m|}$ for $m\ne0$. Counting the
pairs $(a,b)$ with $a-b=\pm m$,
\[
\operatorname{tr}M_D^2=\sum_{a,b}|c_{a-b}|^2=n\alpha^2+2\sum_{m=1}^{n-1}(n-m)\frac{\sin^2\pi m\alpha}{\pi^2m^2}.
\]
Parseval's identity for the indicator of the arc gives $\sum_{m\in\mathbb Z}|c_m|^2=\int_0^{2\pi\alpha}\frac{\dd\theta}{2\pi}=\alpha$, that is $\alpha-\alpha^2=2\sum_{m\ge1}\frac{\sin^2\pi m\alpha}{\pi^2m^2}$. Hence
\begin{align*}
\operatorname{Var}N(D)&=n(\alpha-\alpha^2)-2\sum_{m=1}^{n-1}(n-m)\frac{\sin^2\pi m\alpha}{\pi^2m^2}\\
&=\frac{2n}{\pi^2}\sum_{m\ge1}\frac{\sin^2\pi m\alpha}{m^2}\\
&\quad-\frac2{\pi^2}\sum_{m=1}^{n-1}(n-m)\frac{\sin^2\pi m\alpha}{m^2},
\end{align*}
which is \eqref{eq:numbervariance}. Now let $1\le L\le\frac n2$, so $\frac1n\le\alpha\le\frac12$. In \eqref{eq:numbervariance}, $n\sum_{m\ge n}\frac{\sin^2\pi m\alpha}{m^2}\le n\sum_{m\ge n}\frac1{m^2}\le2$. Split the other sum at $\frac1\alpha$. For $m\le\frac1\alpha$,
$\sin^2\pi m\alpha\le\pi^2m^2\alpha^2$, and these terms add up to at most $\pi^2\alpha^2\sum_{m\le1/\alpha}m\le\pi^2$. For $\frac1\alpha<m<n$ write $\sin^2\pi m\alpha=\frac12-\frac12\cos2\pi m\alpha$. The constant part gives
$\frac12\sum_{1/\alpha<m<n}\frac1m=\frac12\log(n\alpha)+O(1)$. For the oscillating part, the partial sums $S(M)=\sum_{A\le m\le M}\cos2\pi m\alpha$ satisfy $|S(M)|\le\frac2{|1-e^{2\pi i\alpha}|}=\frac1{\sin\pi\alpha}\le\frac1{2\alpha}$
(concavity of $\sin$ on $[0,\frac\pi2]$), and summation by parts gives, with $A$ the first integer above $\frac1\alpha$,
\[
\Bigl|\sum_{m=A}^{B}\frac{\cos2\pi m\alpha}m\Bigr|=\Bigl|\frac{S(B)}B+\sum_{m=A}^{B-1}S(m)\Bigl(\frac1m-\frac1{m+1}\Bigr)\Bigr|\le\frac1{2\alpha}\cdot\frac1A\le\frac12.
\]
Altogether $\sum_{m<n}\frac{\sin^2\pi m\alpha}m=\frac12\log L+O(1)$ with an absolute $O(1)$, and the variance is $\frac1{\pi^2}\log L+O(1)$.
\end{proof}

Compare with $n$ independent uniform points: then $N(D)$ is binomial, with variance $n\alpha(1-\alpha)\approx L$. An arc that expects $1000$ independent points sees about $1000\pm32$ of them; an arc that expects
$1000$ eigenvalues sees $1000\pm0.84$, since $\sqrt{\log1000}/\pi\approx0.84$. The eigenvalues are not random points: they are a crystal that has been gently shaken. Still, a variance that tends to infinity, however slowly, is
enough for a Gaussian limit, and the coin flips make it easy.

\begin{corollary}[Costin--Lebowitz, Soshnikov]\label{cor:countclt}
If $D_n$ are arcs of lengths $2\pi\alpha_n$ with $L_n=n\alpha_n\to\infty$ and $\alpha_n\le\frac12$, then
\[
\frac{N(D_n)-L_n}{\sqrt{\frac1{\pi^2}\log L_n}}\Longrightarrow N(0,1).
\]
\end{corollary}

\begin{proof}
By Proposition~\ref{prop:bernoulli}, $N(D_n)$ has the law of $\sum_jB_{n,j}$ with independent Bernoulli variables of parameters $p_{n,j}$; let $\sigma_n^2=\sum_jp_{n,j}(1-p_{n,j})=\operatorname{Var}N(D_n)$, which tends to infinity
by Theorem~\ref{thm:numbervariance}. The variables $X_{n,j}=(B_{n,j}-p_{n,j})/\sigma_n$ are centred with total variance $1$, and since $|B-p|\le1$,
$\sum_j\mathbb E|X_{n,j}|^3\le\sigma_n^{-3}\sum_j\mathbb E(B_{n,j}-p_{n,j})^2=\sigma_n^{-1}\to0$. Lyapunov's condition holds with $\delta=1$, so $(N(D_n)-L_n)/\sigma_n\Rightarrow N(0,1)$, and $\sigma_n^2\sim\frac1{\pi^2}\log L_n$; conclude with Slutsky's lemma.
\end{proof}

\section{The sine kernel}

On the circle the mean spacing of the eigenvalues is $\frac{2\pi}n$. Measured in this unit, the eigenvalues near the angle $0$ are the points $x_j=\frac n{2\pi}\langle\theta_j\rangle$, where $\langle\theta\rangle\in(-\pi,\pi]$ is the
representative of $\theta$ modulo $2\pi$; on average there is one per unit length. Let
\[
S(u)=\frac{\sin\pi u}{\pi u}\quad(u\ne0),\qquad S(0)=1.
\]

\begin{theorem}[the sine kernel]\label{thm:sinekernel}
For every $k\ge1$ and every bounded Borel function $g$ with compact support on $\R^k$,
\[
\lim_{n\to\infty}\mathbb E\sum_{j_1,\dots,j_k\text{ distinct}}g(x_{j_1},\dots,x_{j_k})=\int_{\R^k}g(x)\det\bigl(S(x_a-x_b)\bigr)_{a,b\le k}\dd x_1\cdots\dd x_k.
\]
\end{theorem}

\begin{proof}
Apply Theorem~8.15 of Volume~8 to the bounded Borel function
\[
\varphi\mapsto g\Bigl(\frac n{2\pi}\langle\varphi_1\rangle,\dots,\frac n{2\pi}\langle\varphi_k\rangle\Bigr),
\]
and substitute $\varphi_a=\frac{2\pi x_a}n$ with $x_a\in(-\frac n2,\frac n2]$, so that
$\frac{\dd\varphi}{(2\pi)^k}=\frac{\dd x}{n^k}$. Summing the geometric series,
\[
K_n(\varphi_a,\varphi_b)=e^{i(n-1)(\varphi_a-\varphi_b)/2}\,\frac{\sin\frac{n(\varphi_a-\varphi_b)}2}{\sin\frac{\varphi_a-\varphi_b}2}.
\]
The phase is $e^{i\psi_a}e^{-i\psi_b}$ with $\psi_a=\frac{(n-1)\varphi_a}2$, so the matrix $(K_n(\varphi_a,\varphi_b))$ is $\Delta A\Delta^{-1}$ with $\Delta=\operatorname{diag}(e^{i\psi_a})$, and has the determinant of $A$. Therefore the expectation equals
\[
\int_{(-\frac n2,\frac n2]^k}g(x)\det\Bigl(S_n(x_a-x_b)\Bigr)_{a,b\le k}\dd x,\qquad S_n(u)=\frac{\sin\pi u}{n\sin\frac{\pi u}n}\quad(S_n(0)=1).
\]
For each $u$, $S_n(u)\to S(u)$, since $n\sin\frac{\pi u}n\to\pi u$. Each $|S_n|\le1$ (as $|K_n|\le n$), so the determinant is bounded by $k!$, and $g$ is bounded with compact support, contained in $(-\frac n2,\frac n2]^k$ for $n$ large.
Dominated convergence gives the limit.
\end{proof}

In words: the rescaled eigenvalues converge to a point process on $\R$ whose $k$-point correlation functions are $\det(S(x_a-x_b))$, the \emph{sine process}. (That correlation functions of this form determine a unique point
process is a theorem of Odile Macchi (1975) and Alexander Soshnikov (2000), which we shall not need: every statement below is about limits of expectations.) For $k=1$ the density is $1$; for $k=2$ it is $1-S(x_1-x_2)^2$,
Dyson's pair correlation.

The determinantal structure is also what makes gaps computable. The matrix $(S(x_a-x_b))$ is positive semidefinite, since $S(u)=\int_{-1/2}^{1/2}e^{2\pi iu\xi}\dd\xi$ makes it a Gram matrix: $\sum_{a,b}\bar v_av_bS(x_a-x_b)=\int_{-1/2}^{1/2}|\sum_bv_be^{2\pi ix_b\xi}|^2\dd\xi$. Likewise
$(K_n(\varphi_a,\varphi_b))$ is the Gram matrix of the vectors $(e^{-ic\varphi_a})_{c<n}$. Both have diagonal entries $1$ and $n$, so by Hadamard's inequality
\begin{equation}\label{eq:detbound}
\begin{gathered}
0\le\det\bigl(S(x_a-x_b)\bigr)_{a,b\le k}\le1,\\
0\le\det\bigl(K_n(\varphi_a,\varphi_b)\bigr)_{a,b\le k}\le n^k.
\end{gathered}
\end{equation}

\begin{theorem}[gap probability]\label{thm:gap}
For $s>0$, the probability that $U$ has no eigenvalue $e^{i\theta}$ with $0\le\theta\le\frac{2\pi s}n$ converges, as $n\to\infty$, to
\[
E(0;s)=1+\sum_{k\ge1}\frac{(-1)^k}{k!}\int_{[0,s]^k}\det\bigl(S(x_a-x_b)\bigr)_{a,b\le k}\dd x_1\cdots\dd x_k.
\]
\end{theorem}

\begin{proof}
Let $D=[0,\frac{2\pi s}n]$ and $n>s$. Expanding the product,
\[
\mathbf 1_{\{N(D)=0\}}=\prod_{j=1}^n\bigl(1-\mathbf 1_D(\theta_j)\bigr)=\sum_{k=0}^n(-1)^k\!\!\sum_{|J|=k}\prod_{j\in J}\mathbf 1_D(\theta_j)=\sum_{k=0}^n\frac{(-1)^k}{k!}\sum_{j_1,\dots,j_k\text{ distinct}}\prod_{a=1}^k\mathbf 1_D(\theta_{j_a}).
\]
Take expectations and apply Volume~8, Theorem~8.15 with $g=\mathbf 1_{D^k}$, then rescale as in the proof of Theorem~\ref{thm:sinekernel}:
\[
P\bigl(N(D)=0\bigr)=\sum_{k=0}^n\frac{(-1)^k}{k!}\int_{[0,s]^k}\det\bigl(S_n(x_a-x_b)\bigr)\dd x.
\]
For $k>n$ the matrix $(K_n(\varphi_a,\varphi_b))_{a,b\le k}$ is a Gram matrix of $k$ vectors in $\mathbb C^n$, so its determinant vanishes and we may let the sum run over all $k\ge0$. By \eqref{eq:detbound} the $k$-th term is at most $\frac{s^k}{k!}$ in
absolute value, uniformly in $n$, and by Theorem~\ref{thm:sinekernel} (with $g=\mathbf 1_{[0,s]^k}$) it converges to the $k$-th term of $E(0;s)$. Dominated convergence for series concludes.
\end{proof}

The series is a \emph{Fredholm determinant}, $E(0;s)=\det(I-\mathcal S_s)$ with $\mathcal S_s$ the integral operator with kernel $S(x-y)$ on $L^2[0,s]$, and that is how one computes it: Folkmar Bornemann showed in 2010 that replacing the integral
by a Gauss--Legendre quadrature with $m$ nodes, $E(0;s)\approx\det\bigl(\delta_{ij}-\sqrt{w_i}S(y_i-y_j)\sqrt{w_j}\bigr)_{i,j\le m}$, converges exponentially fast in $m$. For $n$ independent uniform points the same probability tends to $e^{-s}$.
Near $s=0$, $E(0;s)=1-s+\frac{\pi^2}{36}s^4+O(s^6)$ (Exercise~\ref{ex:smallgap}), against $1-s+\frac{s^2}2-\cdots$ for independent points: the term in $s^2$, the chance of two points close together, has disappeared.

\begin{remark}[spacings]
The gap probability determines the law of the distance between consecutive points: for a stationary point process of intensity $1$, the density of the spacing between a point and the next is $p(s)=\frac{\dd^2}{\dd s^2}E(0;s)$ (a standard fact on stationary point processes, see Mehta's book). For the sine process $p(s)\sim\frac{\pi^2}3s^2$ near $0$, and its tail decays like $e^{-\pi^2s^2/8}$, much faster than the $e^{-s}$ of independent points. Michel Gaudin computed $p$ in 1961, through
prolate spheroidal functions; Wigner had guessed in 1956 a simpler curve, $\frac{\pi s}2e^{-\pi s^2/4}$, for real symmetric matrices, and the analogous guess $\frac{32}{\pi^2}s^2e^{-4s^2/\pi}$ for the unitary case is remarkably close to $p$.
Figure~\ref{fig:zeta} draws $p$, computed by Bornemann's method.
\end{remark}

\section{The characteristic polynomial}

The characteristic polynomial of $U$ on the circle is $Z_n(\theta)=\det(I-e^{-i\theta}U)=\prod_j(1-e^{i(\theta_j-\theta)})$. Since $e^{-i\theta}U$ is again Haar distributed, the law of $Z_n(\theta)$ does not depend on $\theta$, and we write
$Z_n=Z_n(0)=\det(I-U)$. It vanishes exactly at the eigenvalues, and $|Z_n|$ can be as large as $2^n$ (when all eigenvalues are near $-1$). On average it is much more modest.

\begin{proposition}\label{prop:secondmoment}
$\mathbb E|\det(I-U)|^2=n+1$.
\end{proposition}

\begin{proof}
Apply Heine's identity to $f(e^{i\theta})=|1-e^{i\theta}|^2=2-e^{i\theta}-e^{-i\theta}$: $\hat f(0)=2$, $\hat f(\pm1)=-1$, and all other coefficients vanish. The Toeplitz matrix is tridiagonal, with $2$ on the diagonal and $-1$ next to it; expanding
its determinant $D_n$ along the first row gives $D_n=2D_{n-1}-D_{n-2}$, with $D_1=2$ and $D_2=3$. Hence $D_n=n+1$.
\end{proof}

Higher moments by the same route lead to banded Toeplitz determinants that get out of hand quickly. A better idea, due to Paul Bourgade, Chris Hughes, Ashkan Nikeghbali and Marc Yor (2008), is to build $U$ one column at a time.
We need the uniform distribution on the unit sphere of $\mathbb C^k$.

\begin{lemma}[uniform vectors]\label{lem:uniformvector}
Let $g=(g_1,\dots,g_k)$ have independent standard complex Gaussian coordinates (density $\frac1\pi e^{-|z|^2}$ on $\mathbb C$), and $u=g/\|g\|$.
\begin{enumerate}
\item $Vu$ has the law of $u$ for every $V\in U(k)$.
\item If $U$ is Haar distributed on $U(k)$, its first column $Ue_1$ has the law of $u$.
\item $\gamma=u_1$ satisfies $P(|\gamma|^2>x)=(1-x)^{k-1}$ for $0\le x\le1$, so $\mathbb E|\gamma|^{2j}=\frac{j!}{(k)_j}$ with $(k)_j=k(k+1)\cdots(k+j-1)$; and $\mathbb E\bigl[h(|\gamma|)\gamma^a\bar\gamma^b\bigr]=0$ for bounded $h$ and integers $a\ne b$.
\end{enumerate}
\end{lemma}

\begin{proof}
(1) The density of $g$ on $\mathbb C^k=\R^{2k}$ is $\pi^{-k}e^{-\|z\|^2}$, and $V$ is a real orthogonal map of $\R^{2k}$, so $Vg$ has the law of $g$; and $Vu=Vg/\|Vg\|$.

(2) Take $g$ independent of $U$, and let $R_g\in U(k)$ be a unitary matrix depending measurably on $g$ with $R_ge_1=u$ (one is constructed in Lemma~\ref{lem:reflection} below). Given $g$, $UR_g$ is Haar distributed by right invariance, so $UR_ge_1=Uu$ has the
law of $Ue_1$. Given $U$, $Uu$ has the law of $u$ by (1). So $Ue_1$ and $u$ have the same law.

(3) In polar coordinates $P(|g_1|^2>r)=\int_{|z|^2>r}\frac1\pi e^{-|z|^2}\dd z=e^{-r}$: the $E_j=|g_j|^2$ are independent exponential variables, with $\mathbb Ee^{-tE_j}=\frac1{1+t}$. With $T=E_2+\dots+E_k$, independent of $E_1$,
\[
P(|\gamma|^2>x)=P\bigl(E_1(1-x)>xT\bigr)=\mathbb E\,e^{-xT/(1-x)}=\Bigl(1+\frac x{1-x}\Bigr)^{-(k-1)}=(1-x)^{k-1}.
\]
Hence $\mathbb E|\gamma|^{2j}=\int_0^1jx^{j-1}(1-x)^{k-1}\dd x=j\frac{(j-1)!\,(k-1)!}{(j+k-1)!}=\frac{j!}{(k)_j}$. Finally $e^{i\varphi}g_1$ is again standard complex Gaussian and independent of $g_2,\dots,g_k$, so $(|\gamma|,\gamma)$ and
$(|\gamma|,e^{i\varphi}\gamma)$ have the same law; the expectation in question is multiplied by $e^{i(a-b)\varphi}$ for every $\varphi$, so it is $0$.
\end{proof}

For $k=1$, $\gamma$ is uniform on the circle; for $k\ge2$, $|\gamma|<1$ almost surely, and the tail bound gives, for $t\ge0$,
\begin{equation}\label{eq:logtail}
\begin{split}
P\bigl(-\log(1-|\gamma|)>t\bigr)&\le P\bigl(1-|\gamma|^2<2e^{-t}\bigr)\\
&=\min(1,2e^{-t})^{k-1}\le2e^{-t},
\end{split}
\end{equation}
so $\Lambda=-\log(1-|\gamma|)$ has finite moments of all orders.

\begin{lemma}[a reflection]\label{lem:reflection}
Let $m$ be a unit vector of $\mathbb C^k$ with $m\ne e_1$, put $w=e_1-m$, $c=1-\bar m_1$, and $R_m=I-\frac1cww^*$ (and $R_{e_1}=I$). Then $R_m$ is unitary, $R_me_1=m$, and for every $V'\in U(k-1)$
\[
\det\bigl(I-R_m\operatorname{diag}(1,V')\bigr)=(1-m_1)\det(I-V').
\]
\end{lemma}

\begin{proof}
Since $m\ne e_1$ and $|m_1|\le1$, $m_1\ne1$ and $c\ne0$. Now $w^*e_1=1-\bar m_1=c$, so $R_me_1=e_1-w=m$. Next $\|w\|^2=|1-m_1|^2+(1-|m_1|^2)=2-2\operatorname{Re}m_1=2\operatorname{Re}c$, and $\frac1c+\frac1{\bar c}=\frac{2\operatorname{Re}c}{|c|^2}$, so
\[
R_m^*R_m=I-\Bigl(\frac1c+\frac1{\bar c}\Bigr)ww^*+\frac{\|w\|^2}{|c|^2}ww^*=I.
\]
For the determinant, let $D=\operatorname{diag}(1,V')$ and $A=I-D=\operatorname{diag}(0,I-V')$. Then $I-R_mD=A+\frac1cw(D^*w)^*$. For any square $A$ and vectors $x,y$, $\det(A+xy^*)=\det A+y^*\operatorname{adj}(A)x$: for invertible $A$ this is
$\det A\,(1+y^*A^{-1}x)$, and both sides are polynomials in the entries of $A$, which agree on the dense set of invertible matrices. Here $\det A=0$, and $\operatorname{adj}(A)=\det(I-V')\,e_1e_1^{\mathsf T}$: every cofactor of $A$ other than the
$(1,1)$ one contains the zero first row or the zero first column. With $x=w$ and $y=\frac1{\bar c}D^*w$, whose first coordinate is $\frac{w_1}{\bar c}$,
\[
\det(I-R_mD)=\det(I-V')\,\frac{\bar w_1w_1}{c}=\det(I-V')\frac{|1-m_1|^2}{1-\bar m_1}=(1-m_1)\det(I-V').
\]
\end{proof}

\begin{theorem}[Bourgade--Hughes--Nikeghbali--Yor]\label{thm:bhny}
Let $\gamma_1,\dots,\gamma_n$ be independent, $\gamma_k$ distributed as the first coordinate of a uniform vector on the unit sphere of $\mathbb C^k$. If $U$ is Haar distributed on $U(n)$,
\[
\det(I-U)\overset{d}{=}\prod_{k=1}^n(1-\gamma_k).
\]
\end{theorem}

\begin{proof}
By induction on $n$. For $n=1$, $U=e^{i\theta}$ with $\theta$ uniform, and $\gamma_1$ is uniform on the circle. Let $n\ge2$, let $m=Ue_1$ and $W=R_m^{-1}U$. Then $We_1=e_1$, and since $W$ is unitary its first row is also $e_1^{\mathsf T}$, so
$W=\operatorname{diag}(1,U')$ with $U'\in U(n-1)$, and $U=R_m\operatorname{diag}(1,U')$. By Lemma~\ref{lem:reflection}, $\det(I-U)=(1-m_1)\det(I-U')$ (also when $m=e_1$, both sides being $0$).

We claim that $m$ and $U'$ are independent, and that $U'$ is Haar distributed on $U(n-1)$. Fix $V'\in U(n-1)$. The matrix $U\operatorname{diag}(1,V')$ is Haar distributed (right invariance), has the same first column $m$, and
$R_m^{-1}U\operatorname{diag}(1,V')=\operatorname{diag}(1,U'V')$. So $(m,U')$ and $(m,U'V')$ have the same law. For bounded Borel $\phi$ and $\psi$, integrate $\mathbb E[\phi(m)\psi(U'V')]=\mathbb E[\phi(m)\psi(U')]$ against
the Haar measure $\mu$ of $U(n-1)$ in $V'$ and use Fubini: since $\int\psi(xV')\dd\mu(V')=\int\psi\dd\mu$ for every $x$ (left invariance),
\[
\mathbb E[\phi(m)\psi(U')]=\mathbb E\Bigl[\phi(m)\int\psi(U'V')\dd\mu(V')\Bigr]=\mathbb E[\phi(m)]\int\psi\dd\mu.
\]
This is the claim. By Lemma~\ref{lem:uniformvector}(2), $m_1$ has the law of $\gamma_n$, and by the induction hypothesis $\det(I-U')$ has the law of $\prod_{k<n}(1-\gamma_k)$.
\end{proof}

A random unitary matrix is thus a product of $n$ independent reflections, and its characteristic polynomial at $1$ a product of $n$ independent factors, the $k$-th of which is close to $1$ when $k$ is large ($\mathbb E|\gamma_k|^2=\frac1k$). Everything
about $|Z_n|$ is now a matter of independent random variables. Proposition~\ref{prop:secondmoment} is the case $m=1$ of the next theorem, as $\prod_{k=1}^n\frac{k+1}k=n+1$.

\begin{theorem}[Keating--Snaith moments]\label{thm:ksmoments}
For every integer $m\ge0$,
\[
\mathbb E|\det(I-U)|^{2m}=\prod_{l=0}^{m-1}\prod_{i=1}^m\frac{n+l+i}{l+i}\qquad\text{and}\qquad\lim_{n\to\infty}\frac{\mathbb E|\det(I-U)|^{2m}}{n^{m^2}}=\frac{G(m+1)^2}{G(2m+1)},
\]
where $G(m+1)=0!\,1!\cdots(m-1)!$.
\end{theorem}

\begin{proof}
First one factor. Expand
\[
|1-\gamma|^{2m}=(1-\gamma)^m(1-\bar\gamma)^m=\sum_{a,b}\binom ma\binom mb(-\gamma)^a(-\bar\gamma)^b.
\]
By Lemma~\ref{lem:uniformvector}(3), only $a=b$ survives:
\[
\mathbb E|1-\gamma_k|^{2m}=\sum_{j=0}^m\binom mj^2\frac{j!}{(k)_j}.
\]
Multiply by $(k)_m$. Since $\frac{(k)_m}{(k)_j}=(k+j)_{m-j}=(m-j)!\binom{k+m-1}{m-j}$ and $\binom mjj!\,(m-j)!=m!$,
\[
(k)_m\,\mathbb E|1-\gamma_k|^{2m}=m!\sum_{j=0}^m\binom mj\binom{k+m-1}{m-j}=m!\binom{k+2m-1}m=(k+m)_m,
\]
by Vandermonde's identity. So $\mathbb E|1-\gamma_k|^{2m}=\frac{(k+m)_m}{(k)_m}=\prod_{l=0}^{m-1}\frac{k+m+l}{k+l}$. By Theorem~\ref{thm:bhny} and independence,
\[
\mathbb E|Z_n|^{2m}=\prod_{l=0}^{m-1}\prod_{k=1}^n\frac{k+l+m}{k+l}=\prod_{l=0}^{m-1}\frac{(n+l+m)!\,l!}{(l+m)!\,(n+l)!}=\prod_{l=0}^{m-1}\prod_{i=1}^m\frac{n+l+i}{l+i}.
\]
This is a polynomial in $n$ of degree $m^2$ with leading coefficient $\prod_{l=0}^{m-1}\frac{l!}{(l+m)!}=\frac{G(m+1)}{G(2m+1)/G(m+1)}$.
\end{proof}

For $m=2$,
\[
\mathbb E|Z_n|^4=\frac{(n+1)(n+2)^2(n+3)}{12}.
\]
The constants $\frac{G(m+1)^2}{G(2m+1)}$ begin
\[
1,\frac1{12},\frac1{8640},\frac1{870912000},\dots,
\]
which look unremarkable until they are multiplied by $(m^2)!$: the
products are $1,2,42,24024$. Keep these numbers in mind.

\begin{theorem}[a central limit theorem for $\log|\det|$]\label{thm:logclt}
As $n\to\infty$,
\[
\frac{\log|\det(I-U)|}{\sqrt{\frac12\log n}}\Longrightarrow N(0,1).
\]
\end{theorem}

\begin{proof}
By Theorem~\ref{thm:bhny} we may replace $\log|\det(I-U)|$ by $\sum_{k=1}^nX_k$ with $X_k=\log|1-\gamma_k|$ independent (finite almost surely). Fix $k\ge2$. Then $|\gamma_k|<1$, the series $L_k=\sum_{j\ge1}\frac{\gamma_k^j}j$ converges to
$-\log(1-\gamma_k)$, $X_k=-\operatorname{Re}L_k$, and $\sum_j\frac{|\gamma_k|^j}j=\Lambda_k=-\log(1-|\gamma_k|)$ has moments of all orders by \eqref{eq:logtail}. Expanding products of the series and exchanging sums and expectations (justified by
$\mathbb E\Lambda_k^4<\infty$), Lemma~\ref{lem:uniformvector}(3) kills every term in which the powers of $\gamma_k$ and $\bar\gamma_k$ differ. Thus $\mathbb EX_k=0$, $\mathbb EL_k^2=0$, and
\[
\operatorname{Var}X_k=\frac14\mathbb E\bigl(L_k+\bar L_k\bigr)^2=\frac12\mathbb E|L_k|^2=\frac12\sum_{j\ge1}\frac{\mathbb E|\gamma_k|^{2j}}{j^2}=\frac12\sum_{j\ge1}\frac{j!}{j^2(k)_j}.
\]
The term $j=1$ is $\frac1{2k}$. For $j\ge2$, $\frac{j!}{(k)_j}=\prod_{i=1}^j\frac i{k+i-1}\le\frac2{k(k+1)}$, since each factor is at most $1$; so the other terms add up to at most $\frac1{k^2}\sum_{j\ge2}\frac1{j^2}\le\frac1{k^2}$, and $\operatorname{Var}X_k=\frac1{2k}+O(k^{-2})$.

For the fourth moment, $\mathbb EX_k^4\le\mathbb E|L_k|^4=\sum_{s\ge2}\mathbb E|\gamma_k|^{2s}c_s^2$ with $c_s=\sum_{a+b=s}\frac1{ab}=\frac2sH_{s-1}\le2$ ($H$ the harmonic numbers). For $k\ge3$, $\frac{s!}{(k)_s}\le\frac2{k(k+1)}\prod_{i=3}^s\frac i{i+2}=\frac2{k(k+1)}\cdot\frac{12}{(s+1)(s+2)}$, so
\[
\mathbb EX_k^4\le4\cdot\frac{24}{k(k+1)}\sum_{s\ge2}\frac1{(s+1)(s+2)}=\frac{32}{k(k+1)}.
\]
Let $s_n^2=\sum_{k=3}^n\operatorname{Var}X_k=\frac12\log n+O(1)$ and $X_{n,k}=X_k/s_n$ for $3\le k\le n$. These are independent, centred, with total variance $1$, and $\sum_k\mathbb EX_{n,k}^4\le\frac{32}{s_n^4}\sum_k\frac1{k(k+1)}\to0$. Lyapunov's condition
(with $\delta=2$) gives $\sum_{k=3}^nX_k/s_n\Rightarrow N(0,1)$. Finally $(X_1+X_2)/s_n\to0$ almost surely and $s_n/\sqrt{\frac12\log n}\to1$; Slutsky's lemma concludes.
\end{proof}

\begin{think}
The mean of $|\det(I-U)|^2$ is $n+1$, yet $\log|\det(I-U)|$ has mean $0$ (Exercise~\ref{ex:meanlog}). Before reading on, decide which of the two describes a typical matrix. (A typical matrix has $|\det(I-U)|$ of size $e^{O(\sqrt{\log n})}$, so the mean square comes from rare matrices; Direct numerical experimentation shows how hard they are to sample.)
\end{think}

So $|\det(I-U)|$ is typically $e^{\pm\sqrt{\frac12\log n}}$, which is tiny compared with its mean square $n+1$: the moments are dominated by rare matrices whose eigenvalues stay away from $1$. The same argument, with
$\operatorname{Im}L_k$ in place of $\operatorname{Re}L_k$, shows that the imaginary part of $\log\det(I-U)$, defined as $\sum_k\arg(1-\gamma_k)$, also satisfies a central limit theorem with the same normalisation, and that the
two parts are asymptotically independent (Exercise~\ref{ex:argclt}).

\section{The zeros of zeta}

Here is the dictionary. The zeros $\frac12+i\gamma$ of $\zeta$ with $0<\gamma\le T$ number $N(T)=\frac T{2\pi}\log\frac T{2\pi e}+O(\log T)$ (the Riemann--von Mangoldt formula, Volume~5, Chapter~6), so near height $T$ their mean spacing is
$\frac{2\pi}{\log(T/2\pi)}$; the eigenangles of $U\in U(n)$ have mean spacing $\frac{2\pi}n$. The dictionary is $n\leftrightarrow\log\frac T{2\pi}$, and the analogue of $Z_n$ is $\zeta(\frac12+it)$ for $t$ chosen at random near $T$. We first prove one
more statement on the matrix side, the one that Montgomery's theorem mirrors exactly.

\begin{proposition}[traces of powers]\label{prop:formfactor}
For integers $j,l\ge1$, $\mathbb E\operatorname{tr}U^j=0$, $\mathbb E\bigl[\operatorname{tr}U^j\,\overline{\operatorname{tr}U^l}\bigr]=0$ if $j\ne l$, and
\[
\mathbb E|\operatorname{tr}U^j|^2=\min(j,n).
\]
\end{proposition}

\begin{proof}
Replacing $U$ by the Haar distributed $e^{i\varphi}U$ multiplies $\operatorname{tr}U^j\,\overline{\operatorname{tr}U^l}$ by $e^{i(j-l)\varphi}$ and $\operatorname{tr}U^j$ by $e^{ij\varphi}$; so these expectations vanish for $j\ne l$ and $j\ge1$. Next,
$|\operatorname{tr}U^j|^2=n+\sum_{a\ne b}e^{ij(\theta_a-\theta_b)}$, and by Volume~8, Theorem~8.15 with $k=2$ (applied to real and imaginary parts),
\[
\mathbb E\sum_{a\ne b}e^{ij(\theta_a-\theta_b)}=\iint e^{ij(\varphi_1-\varphi_2)}\Bigl(n^2-\sum_{c,d=0}^{n-1}e^{i(c-d)(\varphi_1-\varphi_2)}\Bigr)\frac{\dd\varphi_1\dd\varphi_2}{(2\pi)^2}=0-\#\{(c,d):d-c=j\},
\]
since only the exponentials with total frequency $0$ survive. There are $\max(n-j,0)$ such pairs, and $n-\max(n-j,0)=\min(j,n)$.
\end{proof}

In terms of the rescaled pair correlation: $\frac1n\mathbb E|\operatorname{tr}U^j|^2=\min(\frac jn,1)$, and since $S(u)^2$ is the Fourier transform of the triangle $(1-|\alpha|)_+$, the function $\min(|\alpha|,1)=1-(1-|\alpha|)_+$ is, up to a spike at $\alpha=0$, the Fourier transform of the pair correlation $1-S(u)^2$, read at the frequency $\alpha=\frac jn$. Montgomery's
function is the zeta version. For $T\ge2$ and real $\alpha$ let
\[
F(\alpha)=\frac{2\pi}{T\log T}\sum_{0<\gamma,\gamma'\le T}T^{i\alpha(\gamma-\gamma')}\,\frac4{4+(\gamma-\gamma')^2},
\]
the sum running over pairs of zeros counted with multiplicity; the weight only keeps the sum convergent and does not matter at the scale of the spacing.

\begin{theorem}[Montgomery, 1973; proved in Volume 13]\label{thm:montgomery}
Assume the Riemann hypothesis. For every $\varepsilon>0$, uniformly for $0\le\alpha\le1-\varepsilon$,
\[
F(\alpha)=\bigl(1+o(1)\bigr)T^{-2\alpha}\log T+\alpha+o(1)\qquad(T\to\infty).
\]
\end{theorem}

The proof is an application of the explicit formula of Volume~5, Chapter~7, and Volume~13 gives it in full as Theorem~14.15. Compare with Proposition~\ref{prop:formfactor}: the term $T^{-2\alpha}\log T$ is the trace of $U^0$, a spike at $\alpha=0$, and $\alpha$ is
$\frac jn$. Montgomery conjectured that $F(\alpha)=1+o(1)$ for $\alpha\ge1$, uniformly on bounded intervals, which is $\min(\frac jn,1)=1$ for $j\ge n$, proved above for matrices. By Fourier inversion the conjecture is equivalent to:

\begin{beauty}
\textbf{Montgomery's pair correlation conjecture.} Let $\tilde\gamma=\frac\gamma{2\pi}\log\frac\gamma{2\pi e}$ be the zeros rescaled to mean spacing $1$. For $a<b$, as $T\to\infty$,
\[
\frac1{N(T)}\#\bigl\{(\gamma,\gamma'):0<\gamma\ne\gamma'\le T,\ \tilde\gamma-\tilde\gamma'\in[a,b]\bigr\}\longrightarrow\int_a^b\bigl(1-S(u)^2\bigr)\dd u.
\]
The zeros of zeta, assuming the Riemann hypothesis, would then be indistinguishable from the eigenvalues of a random unitary matrix as far as pairs are concerned, and the natural extension, that all the correlation functions
are the determinants of Theorem~\ref{thm:sinekernel}, was proved by Zeev Rudnick and Peter Sarnak in 1996 for test functions with restricted Fourier support, Montgomery's restriction in higher dimensions. Figure~\ref{fig:zeta}
shows the spacings of $9533$ consecutive zeros above height $10^6$ against the spacing density of the sine process.
\end{beauty}

\begin{figure}[ht]
\centering
\fig{ch11-zeta}
\caption{Left: the spacings between consecutive rescaled zeros $\tilde\gamma$ of $\zeta$ with $10^6<\gamma<10^6+5000$, computed numerically with the Riemann--Siegel formula (all $9533$ zeros predicted by the
Riemann--von Mangoldt formula are found), against the spacing density $p(s)$ of the sine process (red) and the exponential density of independent points (dashed). Right: $\log|\det(I-U)|$ for $20000$ Haar random $U\in U(50)$, divided
by its exact standard deviation $\sqrt{2.74}$, against the standard Gaussian density. The normalisation $\sqrt{\frac12\log50}=\sqrt{1.96}$ of Theorem~\ref{thm:logclt} is only reached slowly, and the left tail, made of the matrices
with an eigenvalue very close to $1$, is still heavy at $n=50$.}
\label{fig:zeta}
\end{figure}

The moments come next. Hardy and Littlewood proved in 1918 that $\frac1T\int_0^T|\zeta(\frac12+it)|^2\dd t\sim\log T$, and Albert Ingham in 1926 that $\frac1T\int_0^T|\zeta(\frac12+it)|^4\dd t\sim\frac1{2\pi^2}(\log T)^4$. No higher moment is known. It was
conjectured that
\[
\frac1T\int_0^T|\zeta(\tfrac12+it)|^{2m}\dd t\sim a_m\frac{g_m}{(m^2)!}(\log T)^{m^2},\qquad a_m=\prod_p\Bigl(1-\frac1p\Bigr)^{m^2}\sum_{r\ge0}\binom{m+r-1}r^2p^{-r},
\]
with an arithmetic factor $a_m$ and a mysterious integer $g_m$: $g_1=1$, $g_2=2$ (Ingham, since $a_2=\frac6{\pi^2}$), and Brian Conrey and Amit Ghosh (1998) and Conrey and Steve Gonek (2001) predicted $g_3=42$ and
$g_4=24024$ by long number-theoretic heuristics. Keating and Snaith's suggestion was to take the arithmetic factor from the primes and the rest from the matrices, with $n=\log\frac T{2\pi}$: by Theorem~\ref{thm:ksmoments},
\[
\frac{g_m}{(m^2)!}=\frac{G(m+1)^2}{G(2m+1)},
\]
which gives $1,2,42,24024$ and every $g_m$ after. Volume~14 tells the rest of the story, together with the analogue of Theorem~\ref{thm:logclt} for zeta, which is older than random matrices: Atle Selberg proved in 1946 that
$\log|\zeta(\frac12+it)|/\sqrt{\frac12\log\log T}$, for $t$ uniform in $[T,2T]$, converges to a standard Gaussian. The dictionary $n\leftrightarrow\log T$ turns $\sqrt{\frac12\log n}$ into Selberg's $\sqrt{\frac12\log\log T}$.

\begin{remark}[the Gaussian unitary ensemble]
Dyson and Montgomery spoke of Hermitian matrices, not unitary ones. The \emph{Gaussian unitary ensemble} is the random Hermitian matrix with density proportional to $e^{-\operatorname{tr}H^2/2}$; its eigenvalues have joint density proportional to
$\prod_{i<j}|\lambda_i-\lambda_j|^2e^{-\sum_i\lambda_i^2/2}$, the same squared Vandermonde as Weyl's formula with a Gaussian weight, and their correlation functions are determinants of a kernel built from Hermite polynomials. In the bulk of the
semicircle, at the scale of the local spacing, this kernel converges to the sine kernel, so the local statistics of the two ensembles coincide (see Anderson, Guionnet and Zeitouni, Chapter~3). Since 2010 the work of
L\'aszl\'o Erd\H os, Horng-Tzer Yau and their collaborators, and of Terence Tao and Van Vu, has shown that the sine kernel appears in the bulk of every Wigner matrix of Chapter~10: local universality, the strongest form of
Wigner's original bet.
\end{remark}

\begin{remark}[general moments]
The decomposition gives every moment, not only the integer ones: for real $t>-\frac12$,
\[
\mathbb E|\det(I-U)|^{2t}=\prod_{k=1}^n\frac{\Gamma(k)\Gamma(k+2t)}{\Gamma(k+t)^2},
\]
the formula in which Keating and Snaith first found their constants, through Selberg's integral (Exercise~\ref{ex:kstreal}). As $n\to\infty$ it behaves like $\frac{G(1+t)^2}{G(1+2t)}n^{t^2}$ with Barnes' $G$-function.
\end{remark}

\section{Exercises}

\begin{exercise}[$\star$]\label{ex:meandet}
Show that $\mathbb E\det(I-U)=1$, in two ways: with Heine's identity, and with Theorem~\ref{thm:bhny}.
\end{exercise}

\begin{exercise}[$\star$]\label{ex:varbound}
Show that $\operatorname{Var}N(D)\le\min\bigl(\mathbb EN(D),\,n-\mathbb EN(D)\bigr)$ for every Borel $D$. Which sets $D$ give variance $0$?
\end{exercise}

\begin{exercise}[$\star\star$]\label{ex:smallgap}
Show that $E(0;s)=1-s+\frac{\pi^2}{36}s^4+O(s^6)$ as $s\to0$, and deduce that the spacing density of the remark on spacings behaves like $\frac{\pi^2}3s^2$ near $0$, assuming it equals $\frac{\dd^2}{\dd s^2}E(0;s)$.
\end{exercise}

\begin{exercise}[$\star\star$]\label{ex:linearstat}
Let $f(e^{i\theta})=\sum_{|j|\le J}\hat f(j)e^{ij\theta}$ be a real trigonometric polynomial, and $\Sigma_f=\sum_{k=1}^nf(e^{i\theta_k})$. Show that
\[
\operatorname{Var}\Sigma_f=\sum_{j\in\mathbb Z}|\hat f(j)|^2\min(|j|,n),
\]
so that for $n\ge J$ the variance no longer depends on $n$. Compare with $n$ independent uniform points.
\end{exercise}

\begin{exercise}[$\star\star$]\label{ex:meanlog}
Show that $\int_0^{2\pi}\log|1-e^{i\theta}|\dd\theta=0$, and deduce that $\mathbb E\log|\det(I-U)|=0$ for every $n$.
\end{exercise}

\begin{exercise}[$\star\star$]\label{ex:argclt}
With the notation of the proof of Theorem~\ref{thm:logclt}, let $Y_k=-\operatorname{Im}L_k$ for $k\ge2$. Show that $\mathbb EY_k^2=\mathbb EX_k^2$ and $\mathbb EX_kY_k=0$, and deduce that $\bigl(\sum_kX_k,\sum_kY_k\bigr)/\sqrt{\frac12\log n}$ converges in law to a
pair of independent standard Gaussians. (Use the Cram\'er--Wold device: a linear combination $aX_k+bY_k$ is again a real part.)
\end{exercise}

\begin{exercise}[$\star\star\star$]\label{ex:ratio}
Let $w=e^{i\varphi}$. Show that $\mathbb E\bigl[\det(I-U)\,\overline{\det(I-wU)}\bigr]=1+\bar w+\dots+\bar w^n$, and deduce that for real $x$
\[
\lim_{n\to\infty}\frac1n\mathbb E\Bigl[Z_n(0)\,\overline{Z_n\bigl(\tfrac{2\pi x}n\bigr)}\Bigr]=e^{i\pi x}\,\frac{\sin\pi x}{\pi x}.
\]
The sine kernel reappears in the characteristic polynomial itself.
\end{exercise}

\begin{exercise}[$\star\star\star$]\label{ex:kstreal}
Let $t>-\frac12$. Show that $\mathbb E|1-\gamma_k|^{2t}=\sum_{j\ge0}\binom tj^2\frac{j!}{(k)_j}$ and, using Gauss's identity $\sum_{j\ge0}\frac{(a)_j(b)_j}{j!\,(c)_j}=\frac{\Gamma(c)\Gamma(c-a-b)}{\Gamma(c-a)\Gamma(c-b)}$ for $c-a-b>0$, deduce the
formula of the remark on general moments. Deduce that $\operatorname{Var}\log|\det(I-U)|=\frac12\sum_{k=1}^n\psi'(k)$, where $\psi'(k)=\sum_{r\ge k}\frac1{r^2}$.
\end{exercise}

\section{Hints}

\begin{hint}{ex:meandet}
With $f(e^{i\theta})=1-e^{i\theta}$ the Toeplitz matrix is triangular. With the decomposition, $\mathbb E\gamma_k=0$.
\end{hint}

\begin{hint}{ex:varbound}
$\sum p_j(1-p_j)$ is at most $\sum p_j$ and at most $\sum(1-p_j)$. Variance $0$ means every $p_j\in\{0,1\}$.
\end{hint}

\begin{hint}{ex:smallgap}
The matrix $(S(x_a-x_b))$ is the Gram matrix of unit vectors $v_a$ (the functions $\xi\mapsto e^{2\pi ix_a\xi}$ in $L^2[-\frac12,\frac12]$), and a Gram determinant is the product of the squared distances of each $v_a$ to the span of the previous ones. Each such distance is at most the distance to $v_{a-1}$ alone. Also $1-S(u)^2\le2(1-S(u))\le\frac{\pi^2}3u^2$.
\end{hint}

\begin{hint}{ex:linearstat}
$\Sigma_f-n\hat f(0)=\sum_{j\ne0}\hat f(j)\operatorname{tr}U^j$, and apply Proposition~\ref{prop:formfactor}.
\end{hint}

\begin{hint}{ex:meanlog}
$\log|1-re^{i\theta}|=-\operatorname{Re}\sum_j\frac{r^je^{ij\theta}}j$ for $r<1$ integrates to $0$; let $r\to1$, dominating $|\log|1-re^{i\theta}||$ by $\log2+|\log|\sin\theta||$.
\end{hint}

\begin{hint}{ex:argclt}
$aX_k+bY_k=-\operatorname{Re}\bigl((a-ib)L_k\bigr)$, and $\mathbb E\bigl((a-ib)L_k\bigr)^2=0$.
\end{hint}

\begin{hint}{ex:ratio}
Heine's identity with $f(e^{i\theta})=(1-e^{i\theta})(1-\bar we^{-i\theta})$ gives a tridiagonal Toeplitz determinant with recursion $D_n=(1+\bar w)D_{n-1}-\bar wD_{n-2}$.
\end{hint}

\begin{hint}{ex:kstreal}
For $|\gamma|<1$, $(1-\gamma)^t=\sum_j\binom tj(-\gamma)^j$ and $|1-\gamma|^{2t}=(1-\gamma)^t(1-\bar\gamma)^t$. Note $\binom tj^2=\frac{(-t)_j^2}{j!^2}$. For the variance, differentiate $\log\mathbb E|Z_n|^{2t}$ twice at $t=0$.
\end{hint}

\section{Solutions}

\begin{solution}{ex:meandet}
\idea The expected determinant only sees the constant term of every expansion, and there is just one.
\writeup Heine: $f(e^{i\theta})=1-e^{i\theta}$ has $\hat f(0)=1$, $\hat f(1)=-1$ and no other coefficient, so the Toeplitz matrix $(\hat f(j-k))$ has $1$ on the diagonal, $-1$ just below it, and $0$ above: it is triangular with determinant
$1$. Decomposition: $\mathbb E\prod(1-\gamma_k)=\prod(1-\mathbb E\gamma_k)=1$ by independence, since $\mathbb E\gamma_k=0$ by Lemma~\ref{lem:uniformvector}(3) with $a=1$, $b=0$.
\end{solution}

\begin{solution}{ex:varbound}
\idea The coin flips do all the work.
\writeup By Proposition~\ref{prop:bernoulli}, $\operatorname{Var}N(D)=\sum_jp_j(1-p_j)$ with $p_j\in[0,1]$. Each term is at most $p_j$ and at most $1-p_j$, and $\sum p_j=\mathbb EN(D)$, $\sum(1-p_j)=n-\mathbb EN(D)$. The variance is $0$ exactly when
every $p_j$ is $0$ or $1$, that is when $\operatorname{tr}M_D^2=\operatorname{tr}M_D$. With $c_m=\int_De^{im\theta}\frac{\dd\theta}{2\pi}$, $\operatorname{tr}M_D^2=\sum_{a=0}^{n-1}\sum_{b=0}^{n-1}|c_{a-b}|^2$, and for each $a$ the inner sum is at most $\sum_{m\in\mathbb Z}|c_m|^2=c_0$ (Parseval), while
$\operatorname{tr}M_D=nc_0$. Equality for $a=0$ forces $c_m=0$ for every $m\ge1$, hence also $c_{-m}=\bar c_m=0$: the indicator of $D$ has no Fourier coefficient but the constant one, so it is constant almost everywhere, $|D|=0$ or $|D|=2\pi$. Then $N(D)$ is almost surely $0$ or $n$. Every other set, however cleverly chosen, has a random count.
\end{solution}

\begin{solution}{ex:smallgap}
\idea Each extra point in a cluster costs a factor $s$ for its position and a factor $s^2$ for its repulsion, so only one and two points matter.
\writeup The term $k=1$ is $-\int_0^s1\dd x=-s$. For $k=2$, $\det=1-S(x_1-x_2)^2$, and $S(u)=1-\frac{\pi^2u^2}6+O(u^4)$ gives $1-S(u)^2=\frac{\pi^2}3u^2+O(u^4)$; since $\iint_{[0,s]^2}(x_1-x_2)^2=\frac{s^4}6$, the term is
$\frac12\bigl(\frac{\pi^2}3\cdot\frac{s^4}6+O(s^6)\bigr)=\frac{\pi^2}{36}s^4+O(s^6)$. For $k\ge3$, write $(S(x_a-x_b))$ as the Gram matrix of the unit vectors $v_a=e^{2\pi ix_a\xi}$ of $L^2[-\frac12,\frac12]$. By Gram--Schmidt, its determinant is $\prod_{a=2}^kd_a^2$ with $d_a$ the
distance from $v_a$ to the span of $v_1,\dots,v_{a-1}$, and $d_a^2\le1-|\langle v_a,v_{a-1}\rangle|^2=1-S(x_a-x_{a-1})^2$. Since $\frac{\sin x}x\ge1-\frac{x^2}6$, $1-S(u)^2\le2(1-S(u))\le\frac{\pi^2}3u^2$. On $[0,s]^k$ the determinant is therefore at most
$(\frac{\pi^2}3s^2)^{k-1}$, and the $k$-th term at most $\frac{s^k}{k!}(\frac{\pi^2}3s^2)^{k-1}$; the terms $k\ge3$ add up to $O(s^7)$. Hence $E(0;s)=1-s+\frac{\pi^2}{36}s^4+O(s^6)$.

For the derivatives: substituting $x=sy$, the $k$-th term is $\frac{(-s)^k}{k!}\int_{[0,1]^k}\det\bigl(S(s(y_a-y_b))\bigr)\dd y$, an entire function of $s$; for complex $|s|\le1$ its entries are bounded by a constant $C$, so the term is at most $|s|^kC^k$ in modulus (expand the
determinant into $k!$ products), and the series converges uniformly for $|s|\le\frac1{2C}$. So $E(0;s)$ is analytic near $0$, its expansion can be differentiated, and $p(s)=E''(0;s)=\frac{\pi^2}3s^2+O(s^4)$. For independent points, $p(s)=e^{-s}$ starts at $1$: small gaps are the commonest of all.
\end{solution}

\begin{solution}{ex:linearstat}
\idea A linear statistic is a combination of traces of powers, and the traces of different powers are uncorrelated.
\writeup Since $\sum_ke^{ij\theta_k}=\operatorname{tr}U^j$, $\Sigma_f=n\hat f(0)+\sum_{0<|j|\le J}\hat f(j)\operatorname{tr}U^j$ (with $\operatorname{tr}U^{-j}=\overline{\operatorname{tr}U^j}$). By Proposition~\ref{prop:formfactor} the traces are centred and
$\mathbb E\bigl[\operatorname{tr}U^j\overline{\operatorname{tr}U^l}\bigr]=\delta_{jl}\min(|j|,n)$ for $j,l\ne0$ (for $j$ and $l$ of opposite signs, $\operatorname{tr}U^j\overline{\operatorname{tr}U^l}$ involves $\operatorname{tr}U^{|j|}\operatorname{tr}U^{|l|}$, which rotation multiplies by
$e^{i(|j|+|l|)\varphi}$, so its expectation also vanishes). Hence $\operatorname{Var}\Sigma_f=\sum_{j\ne0}|\hat f(j)|^2\min(|j|,n)$, which equals $\sum_j|j||\hat f(j)|^2$ as soon as $n\ge J$. For independent uniform points, $\operatorname{Var}\Sigma_f=n\operatorname{Var}f(e^{i\theta})=n\sum_{j\ne0}|\hat f(j)|^2$,
which grows linearly. A smooth statistic of the eigenvalues has bounded fluctuations, without any normalisation: this is the theorem of Diaconis and Shahshahani, and the starting point of a central limit theorem without
$\sqrt n$, due to Kurt Johansson (1997).
\end{solution}

\begin{solution}{ex:meanlog}
\idea Average the logarithm on a smaller circle, where the series converges, and let the circle grow.
\writeup For $r<1$, $\log|1-re^{i\theta}|=-\operatorname{Re}\sum_{j\ge1}\frac{r^je^{ij\theta}}j$, a series converging uniformly in $\theta$; each term integrates to $0$. As $r\to1$, $\log|1-re^{i\theta}|\to\log|1-e^{i\theta}|$ for $\theta\ne0$. For domination: $|1-re^{i\theta}|\le2$, and $|1-re^{i\theta}|\ge|\sin\theta|$, the distance from $1$ to the line through $0$ and $e^{i\theta}$. Since $\log|\sin\theta|$ is
integrable, dominated convergence gives $\int_0^{2\pi}\log|1-e^{i\theta}|\dd\theta=0$, that is $\mathbb EX_1=0$. For $k\ge2$, $\mathbb EX_k=0$ was shown in the proof of Theorem~\ref{thm:logclt}, and all $X_k$ are integrable. So
$\mathbb E\log|\det(I-U)|=\sum_k\mathbb EX_k=0$, although $\log\mathbb E|\det(I-U)|^2=\log(n+1)$: Jensen's inequality is far from equality here.
\end{solution}

\begin{solution}{ex:argclt}
\idea A real part of a rotation-invariant variable has the same law whichever way it is rotated, so the real and imaginary parts carry equal and uncorrelated variance.
\writeup For real $a,b$, $aX_k+bY_k=-\operatorname{Re}\bigl((a-ib)L_k\bigr)$. As in the proof of Theorem~\ref{thm:logclt}, $\mathbb EL_k^2=0$, so $\mathbb E\bigl(\operatorname{Re}(cL_k)\bigr)^2=\frac12|c|^2\mathbb E|L_k|^2$ for every complex $c$. With $c=1$ and $c=-i$:
$\mathbb EX_k^2=\mathbb EY_k^2$; with $c=1-i$: $\mathbb E(X_k+Y_k)^2=2\mathbb EX_k^2$, so $\mathbb EX_kY_k=0$. The fourth moment bound holds for $\operatorname{Re}(cL_k)$ with the factor $|c|^4$. So for every $(a,b)$, Lyapunov's condition gives
$\sum_{k=3}^n(aX_k+bY_k)/s_n\Rightarrow N(0,a^2+b^2)$, which is the law of $aG_1+bG_2$ for independent standard Gaussians; the first two terms are negligible as before. By the Cram\'er--Wold device (Volume~10), the pair converges to $(G_1,G_2)$, and
$s_n\sim\sqrt{\frac12\log n}$.
\end{solution}

\begin{solution}{ex:ratio}
\idea Heine's identity turns the mixed moment into a tridiagonal determinant, which is a geometric series.
\writeup $\det(I-U)\overline{\det(I-wU)}=\prod_jf(e^{i\theta_j})$ with $f(e^{i\theta})=(1-e^{i\theta})(1-\bar we^{-i\theta})=1+\bar w-e^{i\theta}-\bar we^{-i\theta}$, so $\hat f(0)=1+\bar w$, $\hat f(1)=-1$, $\hat f(-1)=-\bar w$. The Toeplitz determinant satisfies
$D_n=(1+\bar w)D_{n-1}-\bar wD_{n-2}$ (expand along the first row; the product of the two off-diagonal entries is $\bar w$), with $D_0=1$, $D_1=1+\bar w$; so $D_n-D_{n-1}=\bar w(D_{n-1}-D_{n-2})=\bar w^n$ and $D_n=\sum_{r=0}^n\bar w^r$.
Now $Z_n(0)\overline{Z_n(\varphi)}=\det(I-U)\overline{\det(I-e^{-i\varphi}U)}$, so $w=e^{-i\varphi}$, $\bar w=e^{i\varphi}$, and for $\varphi=\frac{2\pi x}n$ with $x\ne0$,
\[
\frac1n\sum_{r=0}^ne^{2\pi ixr/n}=\frac1n\cdot\frac{e^{2\pi ix(n+1)/n}-1}{e^{2\pi ix/n}-1}\longrightarrow\frac{e^{2\pi ix}-1}{2\pi ix}=e^{i\pi x}\frac{\sin\pi x}{\pi x},
\]
since $n(e^{2\pi ix/n}-1)\to2\pi ix$; for $x=0$ both sides are $1$ in the limit.
\end{solution}

\begin{solution}{ex:kstreal}
\idea Expand both binomial series, keep the diagonal, and recognise a hypergeometric series at $1$.
\writeup For $k\ge2$ and $|\gamma|<1$, $|1-\gamma|^{2t}=\sum_{a,b}\binom ta\binom tb(-\gamma)^a(-\bar\gamma)^b$. The coefficients satisfy $|\binom ta|\le Ca^{-1-t}$, so $\sum_a|\binom ta|r^a\le C'(1-r)^{-\max(0,-t)}$ for $r<1$ (with a logarithm when $t=0$, a trivial case),
and the square of this is integrable against the law of $r=|\gamma|$, whose density is bounded near $1$ for $k\ge2$, because $2t>-1$. So the expectation may be taken termwise, and Lemma~\ref{lem:uniformvector}(3) keeps $a=b=j$:
$\mathbb E|1-\gamma_k|^{2t}=\sum_j\binom tj^2\frac{j!}{(k)_j}=\sum_j\frac{(-t)_j(-t)_j}{j!\,(k)_j}$. For $k=1$, $\gamma$ is uniform on the circle, $(1-e^{i\theta})^t=\sum_j\binom tj(-1)^je^{ij\theta}$ in $L^2$ (the coefficients are square summable since $2t+2>1$), and
Parseval gives the same formula with $(1)_j=j!$. Gauss's identity with $a=b=-t$, $c=k$ and $c-a-b=k+2t>0$ gives $\frac{\Gamma(k)\Gamma(k+2t)}{\Gamma(k+t)^2}$, and the product over $k$ is the formula of the remark. Then
$\log\mathbb E|Z_n|^{2t}=\sum_k\bigl(\log\Gamma(k)+\log\Gamma(k+2t)-2\log\Gamma(k+t)\bigr)$ has second derivative $\sum_k(4\psi'(k)-2\psi'(k))=2\sum_k\psi'(k)$ at $t=0$, where $\psi=(\log\Gamma)'$. The left side is the cumulant generating function of
$2\log|Z_n|$, whose second derivative at $0$ is $4\operatorname{Var}\log|Z_n|$. So $\operatorname{Var}\log|Z_n|=\frac12\sum_{k=1}^n\psi'(k)$, with $\psi'(k)=\sum_{r\ge k}r^{-2}$ (differentiate $\psi(x+1)=\psi(x)+\frac1x$ and use $\psi'(x)\to0$). Comparing with the proof
of Theorem~\ref{thm:logclt} gives the curious identity $\sum_{j\ge1}\frac{j!}{j^2(k)_j}=\psi'(k)$; for $k=1$ it is $\sum\frac1{j^2}=\frac{\pi^2}6$. For $n=50$ the variance is $2.74$, the value used in Figure~\ref{fig:zeta}.
\end{solution}

\section*{Sources and further reading}

\begin{itemize}
\item \emph{History.} The tea-time meeting of Montgomery and Dyson is told, from interviews with Montgomery, in K.~Sabbagh, \emph{Dr Riemann's Zeros}, Atlantic Books, 2002, pp.~134--136. H.~L.~Montgomery, ``The pair correlation of zeros
of the zeta function'', in \emph{Analytic Number Theory}, Proceedings of Symposia in Pure Mathematics 24, American Mathematical Society, 1973, 181--193. F.~J.~Dyson, ``Statistical theory of the energy levels of complex systems.~I'',
\emph{Journal of Mathematical Physics} 3 (1962), 140--156. M.~Gaudin, ``Sur la loi limite de l'espacement des valeurs propres d'une matrice al\'eatoire'', \emph{Nuclear Physics} 25 (1961), 447--458. A.~M.~Odlyzko, ``On the distribution of
spacings between zeros of the zeta function'', \emph{Mathematics of Computation} 48 (1987), 273--308.
\item \emph{Counts.} J.~B.~Hough, M.~Krishnapur, Y.~Peres and B.~Vir\'ag, ``Determinantal processes and independence'', \emph{Probability Surveys} 3 (2006), 206--229. O.~Costin and J.~L.~Lebowitz, ``Gaussian fluctuation in random
matrices'', \emph{Physical Review Letters} 75 (1995), 69--72. A.~Soshnikov, ``Gaussian fluctuation for the number of particles in Airy, Bessel, sine, and other determinantal random point fields'', \emph{Journal of Statistical Physics} 100
(2000), 491--522. P.~Diaconis and M.~Shahshahani, ``On the eigenvalues of random matrices'', \emph{Journal of Applied Probability} 31A (1994), 49--62. E.~M.~Rains, ``High powers of random elements of compact Lie groups'', \emph{Probability
Theory and Related Fields} 107 (1997), 219--241. F.~Bornemann, ``On the numerical evaluation of Fredholm determinants'', \emph{Mathematics of Computation} 79 (2010), 871--915.
\item \emph{The characteristic polynomial.} J.~P.~Keating and N.~C.~Snaith, ``Random matrix theory and $\zeta(1/2+it)$'', \emph{Communications in Mathematical Physics} 214 (2000), 57--89. P.~Bourgade, C.~P.~Hughes, A.~Nikeghbali and
M.~Yor, ``The characteristic polynomial of a random unitary matrix: a probabilistic approach'', \emph{Duke Mathematical Journal} 145 (2008), 45--69. J.~B.~Conrey and A.~Ghosh, ``A conjecture for the sixth power moment of the Riemann
zeta-function'', \emph{International Mathematics Research Notices} 1998, 775--780. J.~B.~Conrey and S.~M.~Gonek, ``High moments of the Riemann zeta-function'', \emph{Duke Mathematical Journal} 107 (2001), 577--604. Z.~Rudnick and
P.~Sarnak, ``Zeros of principal $L$-functions and random matrix theory'', \emph{Duke Mathematical Journal} 81 (1996), 269--322.
\item M.~L.~Mehta, \emph{Random Matrices}, 3rd ed., Elsevier, 2004. The classical reference for gap probabilities and spacings.
\item E.~S.~Meckes, \emph{The Random Matrix Theory of the Classical Compact Groups}, Cambridge University Press, 2019. Haar matrices, counts, traces of powers and the characteristic polynomial, at exactly the level of this chapter.
\end{itemize}
